% ============================================================
% CAPITOLO 23 - LA V7 NEL BASSO DATO
% ============================================================
\chapter{La \sigla{V}{7}{} nel basso dato}
\label{chap:v7_basso}

Questo capitolo mostra dove la settima di dominante si inserisce in un basso dato, come si risolve fuori dal caso normale, come entra nelle cadenze e come si scrive nella numerica.

\section{Collocazione sulle scale}
\label{sec:v7_basso_dato}

\begin{regola}[Collocazione della \sigla{V}{7}{}]
	\label{reg:collocazione_v7}
	La settima di dominante si costruisce sul \textbf{V grado}, sia del modo maggiore sia del minore (armonico e melodico ascendente). È un accordo \textbf{naturale e indipendente}: \textbf{non richiede la preparazione} della dissonanza (Capitolo \ref{chap:teoria_dissonanza}), perché la sua risoluzione è tanto direzionale da giustificarne l'ingresso «di petto».
\end{regola}

\begin{osservazione}[La settima nel modo minore]
	Nel modo maggiore la settima scende di \textbf{semitono} sulla terza della tonica (in Do maggiore: Fa $\searrow$ Mi); nel modo minore scende di \textbf{tono} (in La minore: Re $\searrow$ Do). Nell'armonizzazione in tonalità minore, ogni volta che la \sigla{V}{7}{} risolve, la terza della tonica di arrivo si raggiunge scendendo di un tono.
\end{osservazione}

\section{Risoluzione d'inganno}
\label{sec:v7_inganno}

Oltre che sulla tonica, la \sigla{V}{7}{} può risolvere sul \textbf{VI grado} (cadenza d'inganno, Capitolo \ref{chap:cadenze}). Le note obbligate del tritono si muovono come nella risoluzione naturale: la sensibile sale alla tonica, la settima scende alla quinta del VI grado.

\begin{regola}[\sigla{V}{7}{} $\rightarrow$ VI]
	\label{reg:v7_inganno}
	La risoluzione d'inganno sul VI grado è ammessa \textbf{purché la triade di arrivo sia completa}. La sensibile sale alla tonica (terza del VI grado), la settima scende di grado sulla quinta del VI.
\end{regola}

\begin{esempio}[\sigla{V}{7}{} completa $\rightarrow$ VI in Do maggiore]
	\label{ex:v7_inganno}
	La sensibile (Si) sale alla tonica (Do); la settima (Fa) scende sulla quinta del VI grado (Mi).
	\musica{es_cap09_04}
\end{esempio}

\section{La \sigla{V}{7}{} nelle cadenze}
\label{sec:v7_cadenze}

La \sigla{V}{7}{} prende il posto della triade di dominante nelle formule cadenzali (Capitolo \ref{chap:cadenze}):

\begin{itemize}
	\item \textbf{cadenza composta consonante}: $\sigla{I}{6}{4} - \sigla{V}{7}{}$;
	\item \textbf{cadenza composta dissonante}: $\sigla{V}{5}{4} - \sigla{V}{7}{3}$;
	\item \textbf{cadenza doppia}: $\sigla{V}{}{} - \sigla{I}{6}{4} - \sigla{V}{5}{4} - \sigla{V}{7}{3}$.
\end{itemize}

\begin{osservazione}[Eccezione di Bach]
	\label{oss:eccezione_bach}
	Se la sensibile si trova in una \textbf{voce interna} e si muove per \textbf{moto contrario} rispetto al basso, può eccezionalmente risolvere \textbf{saltando di terza discendente} sulla quinta della tonica (Si $\searrow$ Sol in Do maggiore), così che la tonica di arrivo sia completa.
\end{osservazione}

Nella cadenza perfetta la settima può diventare essa stessa il ritardo 4-3 sulla tonica (Capitolo \ref{chap:ritardi}): il Fa della \sigla{V}{7}{}, tenuto sul basso Do, forma la quarta che scende alla terza (Mi). Il ritardo può essere realizzato in due forme, come per la risoluzione normale.

\begin{esempio}[Forma A: \sigla{V}{7}{} completa $\rightarrow$ I incompleta con ritardo 4-3 in Do maggiore]
	\label{ex:ritardo_completo_maggiore}
	La \sigla{V}{7}{} usa tutti e quattro i suoni e la tonica è incompleta (senza quinta), con fondamentale raddoppiata e ritardo della terza. Il Fa, nota reale dell'accordo di IV, è tenuto dal contralto e diventa la settima della \sigla{V}{7}{}; nella seconda battuta, sul basso Do, è percosso come ritardo 4 e risolve sul Mi (3).
	\musica{es_cap10_01}
\end{esempio}

\begin{esempio}[Forma B: \sigla{V}{7}{} incompleta $\rightarrow$ I completa con ritardo 4-3 in Do maggiore]
	\label{ex:ritardo_incompleto_maggiore}
	La dominante è senza quinta e con la fondamentale raddoppiata; la settima (Fa) è tenuta e diventa il ritardo. La tonica risultante è completa, con fondamentale, terza (ritardata) e quinta.
	\musica{es_cap10_02}
\end{esempio}

\begin{esempio}[Forma A in La minore]
	\label{ex:ritardo_completo_minore}
	\musica{es_cap10_03}
\end{esempio}

\begin{esempio}[Forma B in La minore]
	\label{ex:ritardo_incompleto_minore}
	\musica{es_cap10_04}
\end{esempio}

\section{La numerica delle settime}
\label{sec:lettura_numerica}

Nel basso figurato la numerica indica gli intervalli da costruire sopra ogni nota del basso. Alle cifre delle triadi (Sezione \ref{subsec:cifratura_accordi}) si aggiungono qui quelle delle settime e le alterazioni.

\begin{definizione}[La numerica del basso figurato]
	\label{def:numerica}
	La \textbf{numerica} (o \textbf{cifratura}) è un sistema di numeri scritti sotto le note del basso che indicano gli intervalli delle voci superiori rispetto al basso stesso. Le convenzioni sono:
	\begin{itemize}
		\item \textbf{nessun numero} (o \textbf{5}): accordo in stato fondamentale;
		\item \textbf{6} (o 6/3): primo rivolto di triade;
		\item \textbf{6/4}: secondo rivolto di triade, da usare solo nei contesti ammessi (Capitolo \ref{chap:quarta_sesta});
		\item \textbf{7}: accordo di settima in stato fondamentale;
		\item \textbf{6/5}: primo rivolto di settima (sesta, quinta e terza sopra il basso);
		\item \textbf{4/3}: secondo rivolto di settima;
		\item \textbf{2} (o 4/2): terzo rivolto di settima, con la settima al basso;
		\item un \textbf{numero preceduto da $\sharp$, $\flat$ o $\natural$}: l'alterazione si applica all'intervallo indicato;
		\item un \textbf{$\sharp$ o $+$ da solo} sotto una nota: si alza di un semitono la terza dell'accordo (frequente nel modo minore, per alzare la sensibile sul V grado).
	\end{itemize}
\end{definizione}

\begin{osservazione}[La numerica è relativa al basso, non alla scala]
	\label{oss:numerica_relativa}
	Gli intervalli della numerica si contano sul \textit{basso reale}, non sulla fondamentale dell'accordo. Se il basso è Re e la numerica è 6, si costruisce una sesta sopra Re, qualunque sia il grado del basso: la nota è quella dell'armatura di chiave (Si in Do maggiore, Si$\flat$ in Fa maggiore). Le alterazioni esplicite nella numerica prevalgono sempre sulle note «naturali» della scala.
\end{osservazione}
